\documentclass[a4paper,12pt]{article}

\usepackage[utf8]{inputenc}
\usepackage[italian]{babel}
\usepackage{amsmath}
\usepackage{geometry}
\usepackage{array}
\usepackage{booktabs}
\usepackage{enumitem}
\geometry{margin=2.0cm}
\usepackage{titlesec}

\titlespacing*{\section}
{0pt}{1ex}{0.5ex}

\titlespacing*{\subsection}
{0pt}{0.8ex}{0.3ex}
\title{Scheda di Laboratorio\\
Verifica delle Leggi di Ohm e di Kirchhoff}
\author{}
\date{}

\begin{document}

\maketitle

\vspace{0.5cm}

\noindent
\textbf{Nome e cognome:} \rule{8cm}{0.4pt}

\vspace{0.3cm}

\noindent
\textbf{Classe:} \rule{3cm}{0.4pt}
\hfill
\textbf{Data:} \rule{3cm}{0.4pt}

\vspace{1cm}

\section*{Obiettivi}

L’esperienza ha lo scopo di:

\begin{itemize}[itemsep=0pt, topsep=1pt]
    \item verificare sperimentalmente la prima legge di Ohm;
    \item verificare le leggi di Kirchhoff;
    \item confrontare i valori teorici con quelli sperimentali;
    \item utilizzare correttamente il multimetro per misure di tensione, corrente e resistenza.
\end{itemize}

\section*{Materiale}

\begin{itemize}[itemsep=0pt, topsep=1pt]
    \item alimentatore in corrente continua;
    \item multimetro digitale;
    \item tre resistenze;
    \item cavetti di collegamento;
    \item breadboard o basetta sperimentale.
\end{itemize}

\section*{Circuito}

Realizzare un circuito costituito da:

\begin{itemize}[itemsep=0pt, topsep=1pt]
    \item una resistenza $R_1$ in serie;
    \item due resistenze $R_2$ e $R_3$ in parallelo.
\end{itemize}

\vspace{0.5cm}

\begin{verbatim}
          R1
 + ----/\/\/\----o----/\/\/\----|
                   |     R2     |
                   |            |
                   ----/\/\/\----
                        R3
                   |
 - ----------------|
\end{verbatim}

\section*{Richiami teorici}

\subsection*{Prima legge di Ohm}

\[
V = RI
\]

\subsection*{Resistenze in parallelo}

\[
\frac{1}{R_p} = \frac{1}{R_2} + \frac{1}{R_3}
\]

\subsection*{Resistenza equivalente}

\[
R_{eq} = R_1 + R_p
\]

\subsection*{Corrente totale}

\[
I_{tot} = \frac{V_{alim}}{R_{eq}}
\]

\subsection*{Prima legge di Kirchhoff}

\[
I_{tot} = I_2 + I_3
\]

\subsection*{Seconda legge di Kirchhoff}

\[
V_{alim} = V_{R1} + V_p
\]

\section*{Parte teorica}

Utilizzando i valori assegnati delle resistenze:

\begin{enumerate}[itemsep=0pt, topsep=1pt]
    \item calcolare la resistenza equivalente del parallelo;
    \item calcolare la resistenza equivalente totale;
    \item calcolare la corrente totale;
    \item calcolare le correnti nei due rami;
    \item calcolare le tensioni ai capi delle resistenze.
\end{enumerate}



\section*{Procedura sperimentale}

\begin{enumerate}[itemsep=0pt, topsep=1pt]
    \item Misurare con il multimetro il valore reale delle resistenze.
    \item Realizzare il circuito assegnato.
    \item Impostare la tensione dell’alimentatore.
    \item Misurare:
    \begin{itemize}[itemsep=0pt, topsep=1pt]
        \item la tensione dell’alimentatore;
        \item la tensione ai capi di $R_1$;
        \item la tensione ai capi del parallelo;
        \item la corrente totale;
        \item la corrente in ciascun ramo.
    \end{itemize}
    \item Riportare i dati nelle tabelle.
    \item Confrontare i dati teorici con quelli sperimentali.
\end{enumerate}

\section*{Errori strumentali}

Assumere come errore strumentale:

\[
\Delta X = \pm \text{ultima cifra significativa}
\]

\vspace{0.5cm}

\section*{Tabella 1 -- Misura delle resistenze}

\begin{center}
\begin{tabular}{|c|c|c|c|}
\hline
Grandezza & Valore nominale & Valore misurato & Errore \\
\hline
$R_1$ & & & \\
\hline
$R_2$ & & & \\
\hline
$R_3$ & & & \\
\hline
\end{tabular}
\end{center}

\vspace{1cm}

\section*{Tabella 2 -- Valori teorici}

\begin{center}
\begin{tabular}{|c|c|}
\hline
Grandezza & Valore teorico \\
\hline
$R_p$ & \\
\hline
$R_{eq}$ & \\
\hline
$I_{tot}$ & \\
\hline
$I_2$ & \\
\hline
$I_3$ & \\
\hline
$V_{R1}$ & \\
\hline
$V_p$ & \\
\hline
\end{tabular}
\end{center}

\vspace{1cm}

\section*{Tabella 3 -- Misure sperimentali}

\begin{center}
\begin{tabular}{|c|c|c|}
\hline
Grandezza & Valore misurato & Errore \\
\hline
$V_{alim}$ & & \\
\hline
$V_{R1}$ & & \\
\hline
$V_p$ & & \\
\hline
$I_{tot}$ & & \\
\hline
$I_2$ & & \\
\hline
$I_3$ & & \\
\hline
\end{tabular}
\end{center}

\vspace{1cm}

\section*{Verifica delle leggi di Kirchhoff}

\subsection*{Legge dei nodi}

\[
I_2 + I_3 = \hspace{4cm}
\]

\[
I_{tot} = \hspace{4cm}
\]

Compatibilità entro gli errori sperimentali:

\vspace{0.3cm}

\noindent
$\fbox{\rule{0pt}{1.2ex}\hspace{1.2ex}}$ Sì \hspace{2cm} $\fbox{\rule{0pt}{1.2ex}\hspace{1.2ex}}$ No

\vspace{1cm}

\subsection*{Legge delle maglie}

\[
V_{R1} + V_p = \hspace{4cm}
\]

\[
V_{alim} = \hspace{4cm}
\]

Compatibilità entro gli errori sperimentali:

\vspace{0.3cm}

\noindent
$\square$ Sì \hspace{2cm} $\square$ No

\vspace{1cm}

\section*{Conclusioni}

Rispondere alle seguenti domande:

\begin{enumerate}[itemsep=0pt, topsep=1pt]
    \item I valori sperimentali sono compatibili con quelli teorici?
    \item Le leggi di Ohm e di Kirchhoff risultano verificate?
    \item Quali sono le principali cause di errore sperimentale?
    \item Le resistenze reali coincidono con i valori nominali?
\end{enumerate}

\vspace{2cm}

\noindent
\textbf{Osservazioni finali}

\vspace{4cm}

\end{document}